
\section*{Science : Le modèle DP (Ullman) et les seuils lexicaux (Nation) comme outils critiques}

Or, ce régime pédagogique repose sur des présupposés que la science cognitive contemporaine invalide. Le cerveau n'est pas un muscle généraliste dont l'exercice sur une langue ancienne transfert automatiquement sa rigueur aux mathématiques ou à la philosophie. La mémoire ne fonctionne pas selon le modèle de la Psychologie des Facultés du XIXe siècle.

Les neurosciences, en particulier le \textbf{modèle Déclaratif/Procédural} de Michael Ullman, montrent que l'apprentissage des langues engage deux systèmes de mémoire distincts : un système déclaratif (les règles explicites, les faits) et un système procédural (les automatismes, les réflexes moteurs). Or, l'enseignement français du grec sature le système déclaratif---il fait mémoriser des paradigmes, des conjugaisons, des règles---tout en négligeant complètement l'entraînement du système procédural.

Parallèlement, les travaux de Paul Nation et Batia Laufer en psycholinguistique établissent un fait décisif : \textbf{on ne peut lire un texte que si l'on connaît au moins 95-98\% des mots}. Or, l'approche française, en insistant sur le déchiffrement grammatical plutôt que sur l'accumulation lexicale massive, rend quasiment impossible l'atteinte de ce seuil. Les élèves restent prisonniers du dictionnaire, jamais libérés pour la lecture vivante et fluide.

Ces deux découvertes scientifiques---l'inadéquation de l'approche explicite grammaticale et l'impératif du seuil lexical---fournissent les outils critiques pour déconstruire le régime pédagogique hérité de 1880.
